\documentclass[../../main.tex]{subfiles}
\begin{document}
% SOURCE p=174; printed=159; object=ch05-linear-functionals-exercises-title; region=lower
\subsection*{习题}
% SOURCE p=174; printed=159; object=ch05-ex17-functional-closed-kernel; number=17; region=lower
\begin{exercise}\label{exercise:ch05-functional-closed-kernel}
赋范向量空间 $\mathcal X$ 上的线性泛函 $f$ 有界，当且仅当 $f^{-1}(\{0\})$ 闭。（使用\subcref{exercise:ch05-quotient-norm}{item:ch05-quotient-almost-unit}。）
\end{exercise}
% SUPPLEMENT: id=v14-ch05-ex17-solution; parent=ch05-ex17-functional-closed-kernel; kind=solution+motivation+hypotheses; origin=independent-derivation
\begin{solution}
\textbf{补写解答（重构新增）。}若 $f$ 有界，则连续，故闭集 $\{0\}$ 的原像闭。反过来，设 $\mathcal M=\ker f$ 闭。若 $f=0$ 则有界。否则选 $z$ 使 $f(z)=1$，则 $z\notin\mathcal M$，距离 $d=\inf_{m\in\mathcal M}\|z-m\|$ 为正。对任意 $x$，有 $x-f(x)z\in\mathcal M$。因此在商空间中 $x+\mathcal M=f(x)(z+\mathcal M)$，由商范数的齐次性和收缩估计，
\begin{align*}
|f(x)|d=\|x+\mathcal M\|\le\|x\|.
\end{align*}
故 $|f(x)|\le d^{-1}\|x\|$，$f$ 有界。这使用了提示所依据的同一个事实：闭子空间外的非零商方向具有正范数；不必要求到子空间的距离真正取到。
\end{solution}
\begin{remark}
\textbf{动机与条件（重构新增）。}标量值线性映射的非零商空间只有一维，因而核闭足以控制整个映射。闭性使 $d>0$；线性性使所有商类由 $z$ 的一个标量倍数描述。这个判别不需要 $\mathcal X$ 完备，也不应直接推广为任意无限维值域算子的有界性判别。
\end{remark}
% END SUPPLEMENT: id=v14-ch05-ex17-solution
% SOURCE p=174; printed=159; object=ch05-ex18-finite-subspace-closed; number=18; region=bottom
\begin{exercise}\label{exercise:ch05-finite-subspace-closed}
设 $\mathcal X$ 为赋范向量空间。
\begin{enumerate}
\item\label{item:ch05-adjoin-one-closed} 若 $\mathcal M$ 为闭子空间且 $x\in\mathcal X\setminus\mathcal M$，则 $\mathcal M+\mathbb Cx$ 闭。（使用\subcref{theorem:ch05-dual-separation-embedding}{item:ch05-dual-separates-subspace}。）
\item\label{item:ch05-finite-dim-closed} $\mathcal X$ 的每个有限维子空间都闭。
\end{enumerate}
\end{exercise}
% SUPPLEMENT: id=v14-ch05-ex18-solution; parent=ch05-ex18-finite-subspace-closed; kind=solution+motivation+hypotheses; origin=independent-derivation
\begin{solution}
\textbf{逐项补写（重构新增）。}
\begin{enumerate}
\item 对应\subcref{exercise:ch05-finite-subspace-closed}{item:ch05-adjoin-one-closed}：由分离定理选连续泛函 $f$，使 $f|_{\mathcal M}=0$ 且 $f(x)=\delta>0$。设 $z_n=y_n+\lambda_nx\to z$，其中 $y_n\in\mathcal M$。连续性给 $\lambda_n=f(z_n)/\delta\to f(z)/\delta=\lambda$，于是 $y_n=z_n-\lambda_nx\to z-\lambda x$。因 $\mathcal M$ 闭，$z-\lambda x\in\mathcal M$，故 $z\in\mathcal M+\mathbb Cx$。赋范空间为度量空间，序列闭性即闭性。
\item 对应\subcref{exercise:ch05-finite-subspace-closed}{item:ch05-finite-dim-closed}：零子空间闭。对有限维子空间选有限基 $e_1,\ldots,e_n$，从零子空间起，逐次添加 $\mathbb Ce_j$。每一步新基向量不属于先前张成的空间，故上一问保持闭性；有限步后得到整个子空间。实空间把 $\mathbb C$ 换为 $\mathbb R$，证明不变。
\end{enumerate}
\end{solution}
\begin{remark}
\textbf{动机与条件（重构新增）。}连续的分离泛函读出新增方向的系数，防止两个分量各自发散而其和收敛。$x\notin\mathcal M$ 用来取得非零读数；若 $x\in\mathcal M$，和空间本来就等于 $\mathcal M$。有限维条件使归纳在有限步结束；无限维子空间未必闭，例如有限支集序列在 $\ell^1$ 中稠密而非闭。环境空间不必完备。
\end{remark}
% END SUPPLEMENT: id=v14-ch05-ex18-solution
% SOURCE p=175; printed=160; object=ch05-ex19-infinite-dimension-noncompact; number=19; region=top
\begin{exercise}\label{exercise:ch05-infinite-dimension-noncompact}
设 $\mathcal X$ 为无限维赋范向量空间。
\begin{enumerate}
\item\label{item:ch05-separated-unit-sequence} 存在序列 $\{x_j\}$，使每个 $j$ 都有 $\|x_j\|=1$，且 $j\ne k$ 时 $\|x_j-x_k\|\ge\frac12$。（使用\subcref{exercise:ch05-quotient-norm}{item:ch05-quotient-almost-unit}和\cref{exercise:ch05-finite-subspace-closed}归纳构造 $x_j$。）
\item\label{item:ch05-no-local-compactness} $\mathcal X$ 不局部紧。
\end{enumerate}
\end{exercise}
% SUPPLEMENT: id=v14-ch05-ex19-solution; parent=ch05-ex19-infinite-dimension-noncompact; kind=solution+motivation+hypotheses; origin=independent-derivation
\begin{solution}
\textbf{逐项补写（重构新增）。}
\begin{enumerate}
\item 对应\subcref{exercise:ch05-infinite-dimension-noncompact}{item:ch05-separated-unit-sequence}：先取单位向量 $x_1$。若已选 $x_1,\ldots,x_n$，则 $\mathcal M_n=\operatorname{span}\{x_1,\ldots,x_n\}$ 有限维，故闭；无限维假设保证它是真子空间。由提示的商范数结论取单位向量 $x_{n+1}$，使 $\inf_{y\in\mathcal M_n}\|x_{n+1}-y\|\ge1/2$。每个先前 $x_j\in\mathcal M_n$，所以与新向量的距离至少 $1/2$。归纳完成。
\item 对应\subcref{exercise:ch05-infinite-dimension-noncompact}{item:ch05-no-local-compactness}：若零点有紧邻域 $K$，可取 $r>0$ 使 $B(0,r)\subseteq K$。令 $z_j=(r/2)x_j$，则 $z_j\in K$，相异两项距离至少 $r/4$。但 $K$ 被所有半径 $r/10$ 的开球覆盖，紧性给有限子覆盖；每个这样的球内至多有一个 $z_j$，因为同球两点距离小于 $r/5<r/4$。这不可能覆盖无穷多个 $z_j$，矛盾。
\end{enumerate}
\end{solution}
\begin{remark}
\textbf{动机与条件（重构新增）。}无限多个新方向使单位球中不断出现与已有有限维部分保持固定距离的点，从而破坏紧性所需的有限小球覆盖。无限维性用于每一步保持真子空间，赋范结构使这种分离能缩放进任意零邻域；不需要完备性。有限维空间则通过范数等价继承欧氏空间的局部紧性。
\end{remark}
% END SUPPLEMENT: id=v14-ch05-ex19-solution
% SOURCE p=175; printed=160; object=ch05-ex20-finite-dimensional-complement; number=20; region=upper
\begin{exercise}\label{exercise:ch05-finite-dimensional-complement}
若 $\mathcal M$ 为赋范向量空间 $\mathcal X$ 的有限维子空间，则存在闭子空间 $\mathcal N$，使 $\mathcal M\cap\mathcal N=\{0\}$ 且 $\mathcal M+\mathcal N=\mathcal X$。
\end{exercise}
% SUPPLEMENT: id=v14-ch05-ex20-solution; parent=ch05-ex20-finite-dimensional-complement; kind=solution+motivation+hypotheses; origin=independent-derivation
\begin{solution}
\textbf{补写解答（重构新增）。}若 $\mathcal M=\{0\}$，取 $\mathcal N=\mathcal X$。否则取 $\mathcal M$ 的基 $e_1,\ldots,e_n$，坐标泛函 $u_j(\sum_k a_ke_k)=a_j$ 在有限维空间 $\mathcal M$ 上有界。用保持范数的 Hahn--Banach 延拓，得到 $f_j\in\mathcal X^*$ 且 $f_j|_{\mathcal M}=u_j$。定义
\begin{align*}
Px=\sum_{j=1}^nf_j(x)e_j.
\end{align*}
它线性，且 $\|Px\|\le\sum_j\|f_j\|\|e_j\|\|x\|$，所以有界。对 $m\in\mathcal M$ 有 $Pm=m$，且所有 $Px\in\mathcal M$，故 $P^2=P$。令 $\mathcal N=\ker P$，它是连续映射下零点的原像，因而闭。若 $x\in\mathcal M\cap\mathcal N$，则 $x=Px=0$。任意 $x$ 都可写成 $Px+(x-Px)$，前项在 $\mathcal M$，后项因 $P(x-Px)=0$ 在 $\mathcal N$，故两空间之和为 $\mathcal X$。
\end{solution}
\begin{remark}
\textbf{动机与条件（重构新增）。}逐个延拓有限多个坐标，便可在全空间上连续地提取有限维分量；核就是剩余分量。有限维性保证所有坐标有界、只有有限项相加。闭补空间比纯代数补空间更强，因为分解由连续投影实现；本题不要求环境空间完备，也不证明每个无限维闭子空间都有闭补空间。
\end{remark}
% END SUPPLEMENT: id=v14-ch05-ex20-solution
% SOURCE p=175; printed=160; object=ch05-ex21-dual-product; number=21; region=upper-middle
\begin{exercise}\label{exercise:ch05-dual-product}
设 $\mathcal X,\mathcal Y$ 为赋范向量空间，定义 $\alpha:\mathcal X^*\times\mathcal Y^*\to(\mathcal X\times\mathcal Y)^*$ 为 $\alpha(f,g)(x,y)=f(x)+g(y)$。则 $\alpha$ 是同构；若在 $\mathcal X\times\mathcal Y$ 上取范数 $\|(x,y)\|=\max(\|x\|,\|y\|)$，在其对偶上取相应算子范数，在 $\mathcal X^*\times\mathcal Y^*$ 上取 $\|(f,g)\|=\|f\|+\|g\|$，则此同构等距。
\end{exercise}
% SUPPLEMENT: id=v14-ch05-ex21-solution; parent=ch05-ex21-dual-product; kind=solution+well-definedness+motivation+hypotheses; origin=independent-derivation
\begin{solution}
\textbf{补写解答（重构新增）。}按题设范数，有
\begin{align*}
|f(x)+g(y)|\le(\|f\|+\|g\|)\max(\|x\|,\|y\|),
\end{align*}
所以像确为有界泛函，$\alpha$ 线性且 $\|\alpha(f,g)\|\le\|f\|+\|g\|$。给定 $\varepsilon>0$，从闭单位球中选 $x,y$，使 $|f(x)|>\|f\|-\varepsilon$、$|g(y)|>\|g\|-\varepsilon$。分别乘单位相位，使两个值成为非负实数，仍记为 $x,y$；零值无需改变。于是
\begin{align*}
\|\alpha(f,g)\|\ge|f(x)+g(y)|>\|f\|+\|g\|-2\varepsilon.
\end{align*}
让 $\varepsilon\downarrow0$ 即得等号。使用闭单位球也涵盖零空间和零泛函，不假定范数上确界能取到。

若 $h\in(\mathcal X\times\mathcal Y)^*$，定义 $f(x)=h(x,0)$、$g(y)=h(0,y)$。两者线性且各自范数至多 $\|h\|$，所以属于对应的对偶；线性性给 $h(x,y)=f(x)+g(y)$，故 $\alpha$ 满射。若 $\alpha(f,g)=0$，分别取 $(x,0)$ 和 $(0,y)$ 得 $f=g=0$，故单射。等距双射的逆也等距，因而是赋范空间同构。换为先前介绍的其他有限乘积范数，它们等价，所以仍为有界且逆有界的同构。
\end{solution}
\begin{remark}
\textbf{动机与条件（重构新增）。}乘积上的线性测量可分成各坐标的测量。最大范数允许两个坐标同时取到各自的单位尺度，所以对偶大小是两项之和；独立相位避免复数值相互抵消。若改变乘积范数，代数分解仍在，但精确的等距公式随之改变。完备性不是此结论的条件。
\end{remark}
% END SUPPLEMENT: id=v14-ch05-ex21-solution
% SOURCE p=175; printed=160; object=ch05-ex22-adjoint; number=22; region=middle
\begin{exercise}\label{exercise:ch05-adjoint}
设 $\mathcal X,\mathcal Y$ 为赋范向量空间，$T\in L(\mathcal X,\mathcal Y)$。
\begin{enumerate}
\item\label{item:ch05-adjoint-norm}
\begin{definition*}[伴随算子]\phantomsection\label{definition:ch05-adjoint-operator}
定义 $T^\dagger:\mathcal Y^*\to\mathcal X^*$ 为 $T^\dagger f=f\circ T$，称 $T^\dagger$ 为 $T$ 的\textbf{伴随}或\textbf{转置}。
\end{definition*}
则 $T^\dagger\in L(\mathcal Y^*,\mathcal X^*)$ 且 $\|T^\dagger\|=\|T\|$。
\item\label{item:ch05-second-adjoint} 将上一问构造应用两次，得到 $T^{\dagger\dagger}\in L(\mathcal X^{**},\mathcal Y^{**})$。若将 $\mathcal X,\mathcal Y$ 分别与其在双对偶中的自然像 $\widehat{\mathcal X},\widehat{\mathcal Y}$ 等同，则 $T^{\dagger\dagger}|_{\mathcal X}=T$。
\item\label{item:ch05-adjoint-injective} $T^\dagger$ 单射，当且仅当 $T$ 的值域在 $\mathcal Y$ 中稠密。
\item\label{item:ch05-adjoint-dense-range} 若 $T^\dagger$ 的值域在 $\mathcal X^*$ 中稠密，则 $T$ 单射；当 $\mathcal X$ 自反时，逆命题也成立。
\end{enumerate}
\end{exercise}
% SUPPLEMENT: id=v14-ch05-ex22-solution; parent=ch05-ex22-adjoint; kind=solution+well-definedness+motivation+hypotheses; origin=independent-derivation
\begin{solution}
\textbf{逐项补写（重构新增）。}
\begin{enumerate}
\item 对应\subcref{exercise:ch05-adjoint}{item:ch05-adjoint-norm}：复合 $f\circ T$ 线性且范数至多 $\|f\|\|T\|$，故定义良好；对 $f$ 的线性组合，复合逐项分配，所以 $T^\dagger$ 线性且 $\|T^\dagger\|\le\|T\|$。对任意 $x$，若 $Tx\ne0$，由范数支撑泛函取 $f\in\mathcal Y^*$，使 $\|f\|=1$、$f(Tx)=\|Tx\|$。于是 $\|Tx\|=|(T^\dagger f)(x)|\le\|T^\dagger\|\|x\|$。若 $Tx=0$ 此界也成立；对所有 $x$ 取最佳界得 $\|T\|\le\|T^\dagger\|$。两个方向合起来得等号，包括 $T=0$ 的情形。
\item 对应\subcref{exercise:ch05-adjoint}{item:ch05-second-adjoint}：第一问给二次伴随有界。对 $x\in\mathcal X$、$f\in\mathcal Y^*$，
\begin{align*}
(T^{\dagger\dagger}\widehat x)(f)=\widehat x(T^\dagger f)=(T^\dagger f)(x)=f(Tx)=\widehat{Tx}(f).
\end{align*}
两个双对偶元素在每个 $f$ 上相等，故 $T^{\dagger\dagger}\widehat x=\widehat{Tx}$，这正是题设等同下的限制关系。
\item 对应\subcref{exercise:ch05-adjoint}{item:ch05-adjoint-injective}：若值域稠密且 $T^\dagger f=0$，则 $f$ 在 $T\mathcal X$ 上为零，由连续性在其闭包 $\mathcal Y$ 上也为零，故核为零。反过来，若闭子空间 $\overline{T\mathcal X}\ne\mathcal Y$，Hahn--Banach 分离给非零 $f\in\mathcal Y^*$ 在该子空间上消失；此时 $T^\dagger f=0$，与单射矛盾。
\item 对应\subcref{exercise:ch05-adjoint}{item:ch05-adjoint-dense-range}：若 $Tx=0$，每个 $f\in\mathcal Y^*$ 都满足 $(T^\dagger f)(x)=0$。评价泛函 $\widehat x$ 连续，所以若 $T^\dagger\mathcal Y^*$ 稠密，它就在整个 $\mathcal X^*$ 上为零。由对偶分离点得 $x=0$，所以 $T$ 单射。

反过来，设 $\mathcal X$ 自反且 $T$ 单射。若 $\mathcal R=\overline{T^\dagger\mathcal Y^*}$ 是 $\mathcal X^*$ 的真闭子空间，分离定理给非零 $F\in\mathcal X^{**}$ 在 $\mathcal R$ 上为零。自反性使 $F=\widehat x$，且由等距性 $x\ne0$。于是对每个 $f\in\mathcal Y^*$，$f(Tx)=\widehat x(T^\dagger f)=0$；对偶分离点给 $Tx=0$，与单射矛盾。所以伴随值域稠密。
\end{enumerate}
\end{solution}
\begin{remark}
\textbf{动机与逐条件说明（重构新增）。}伴随把值域上的测量拉回定义域，因而方向反转；二次伴随再回到原方向，且保持自然嵌入。有界性保证拉回的测量仍连续，Hahn--Banach 的范数支撑保证范数不损失。这里是 Banach 空间意义的伴随，无需内积，复标量线性中没有额外共轭。稠密性使连续泛函由值域上的值决定；最后逆向中，自反性恰好把分离伴随值域的双对偶元素变成真实的 $x\in\mathcal X$。其余三问无需原空间完备。
\end{remark}
% END SUPPLEMENT: id=v14-ch05-ex22-solution
\Needspace{14\baselineskip}
% SOURCE p=175; printed=160; object=ch05-ex23-annihilators; number=23; region=middle-lower
\begin{exercise}\label{exercise:ch05-annihilators}
设 $\mathcal X$ 为 Banach 空间，$\mathcal M$ 为 $\mathcal X$ 的闭子空间，$\mathcal N$ 为 $\mathcal X^*$ 的闭子空间。
\begin{definition*}[本题的零化子]\phantomsection\label{definition:ch05-annihilators}
令 $\mathcal M^0=\{f\in\mathcal X^*:f|_{\mathcal M}=0\}$，并令 $\mathcal N^\perp=\{x\in\mathcal X:f(x)=0\text{ 对所有 }f\in\mathcal N\}$。因此，把 $\mathcal X$ 与其在 $\mathcal X^{**}$ 中的像等同时，有 $\mathcal N^\perp=\mathcal N^0\cap\mathcal X$。
\end{definition*}
\begin{enumerate}
\item\label{item:ch05-annihilators-closed} $\mathcal M^0,\mathcal N^\perp$ 分别为 $\mathcal X^*,\mathcal X$ 的闭子空间。
\item\label{item:ch05-double-annihilator} $(\mathcal M^0)^\perp=\mathcal M$，且 $(\mathcal N^\perp)^0\supseteq\mathcal N$。若 $\mathcal X$ 自反，则 $(\mathcal N^\perp)^0=\mathcal N$。
\item\label{item:ch05-quotient-dual-annihilator} 令 $\pi:\mathcal X\to\mathcal X/\mathcal M$ 为自然投影，定义 $\alpha:(\mathcal X/\mathcal M)^*\to\mathcal X^*$ 为 $\alpha(f)=f\circ\pi$。则 $\alpha$ 是从 $(\mathcal X/\mathcal M)^*$ 到 $\mathcal M^0$ 的等距同构，其中商空间赋以商范数。
\item\label{item:ch05-subspace-dual-quotient} 定义 $\beta:\mathcal X^*\to\mathcal M^*$ 为 $\beta(f)=f|_{\mathcal M}$；它像\cref{exercise:ch05-kernel-factorization}那样诱导映射 $\overline\beta:\mathcal X^*/\mathcal M^0\to\mathcal M^*$，且 $\overline\beta$ 是等距同构。
\end{enumerate}
\end{exercise}
% SUPPLEMENT: id=v14-ch05-ex23-solution; parent=ch05-ex23-annihilators; kind=solution+well-definedness+motivation+hypotheses; origin=independent-derivation
\begin{solution}
\textbf{逐项补写（重构新增）。}
\begin{enumerate}
\item 对应\subcref{exercise:ch05-annihilators}{item:ch05-annihilators-closed}：有 $\mathcal M^0=\bigcap_{x\in\mathcal M}\ker\widehat x$、$\mathcal N^\perp=\bigcap_{f\in\mathcal N}\ker f$。各核均为闭子空间，任意交仍为闭子空间，所以结论成立。
\item 对应\subcref{exercise:ch05-annihilators}{item:ch05-double-annihilator}：每个 $x\in\mathcal M$ 被所有 $\mathcal M^0$ 中的泛函消去，故 $\mathcal M\subseteq(\mathcal M^0)^\perp$。若 $x\notin\mathcal M$，闭子空间分离定理给 $f\in\mathcal M^0$ 且 $f(x)\ne0$，所以 $x\notin(\mathcal M^0)^\perp$，得等号。对于 $f\in\mathcal N$，定义直接给其在 $\mathcal N^\perp$ 上消失，故 $\mathcal N\subseteq(\mathcal N^\perp)^0$。

若 $\mathcal X$ 自反且 $g\notin\mathcal N$，在赋范空间 $\mathcal X^*$ 中对闭子空间 $\mathcal N$ 分离，得到 $F\in\mathcal X^{**}$，使 $F|_{\mathcal N}=0$、$F(g)\ne0$。写 $F=\widehat x$，则 $x\in\mathcal N^\perp$ 而 $g(x)\ne0$，所以 $g\notin(\mathcal N^\perp)^0$。这给反包含。
\item 对应\subcref{exercise:ch05-annihilators}{item:ch05-quotient-dual-annihilator}：$\alpha$ 线性且其值消去 $\mathcal M$，又 $\|\alpha f\|\le\|f\|$。若 $\alpha f=0$，由 $\pi$ 满射得 $f=0$。若 $h\in\mathcal M^0$，令 $f(x+\mathcal M)=h(x)$，换代表时 $h$ 不变，故良定义、线性。对每个 $m\in\mathcal M$，
\begin{align*}
|f(x+\mathcal M)|=|h(x+m)|\le\|h\|\|x+m\|.
\end{align*}
取下确界得 $|f(x+\mathcal M)|\le\|h\|\|x+\mathcal M\|$，故 $f$ 有界且 $\alpha f=h$，证明满射。还得 $\|f\|\le\|h\|=\|\alpha f\|$，与相反界合得等距。零商时公式同样成立。
\item 对应\subcref{exercise:ch05-annihilators}{item:ch05-subspace-dual-quotient}：限制映射 $\beta$ 线性、有界，范数至多1，核恰为 $\mathcal M^0$。Hahn--Banach 保范数延拓给 $\beta$ 满射。故 $\overline\beta(f+\mathcal M^0)=f|_{\mathcal M}$ 良定义、线性、双射。每个代表 $f+h$，$h\in\mathcal M^0$，都限制为 $f|_{\mathcal M}$，因此
\begin{align*}
\|f|_{\mathcal M}\|\le\inf_{h\in\mathcal M^0}\|f+h\|=\|f+\mathcal M^0\|.
\end{align*}
另一方面，取 $f|_{\mathcal M}$ 的保范数延拓 $F$，则 $F-f\in\mathcal M^0$，所以 $\|f+\mathcal M^0\|\le\|F\|=\|f|_{\mathcal M}\|$。合起来给等距，进而逆映射也有界。
\end{enumerate}
\end{solution}
\begin{remark}
\textbf{动机与逐条件说明（重构新增）。}零化子收集看不见某个子空间的测量。取商会让消去该子空间的测量自然下降，限制到子空间则把相差一个零化子元素的测量视为相同，因而得到两种对偶识别。第一问连原集合闭性也不需要；第二问恢复 $\mathcal M$ 本身需要它闭，否则恢复的是闭包。$\mathcal N$ 闭且 $\mathcal X$ 自反，使双对偶分离泛函能由原空间元素实现。第三问的闭性保证商范数正定；第四问的精确范数来自保范数延拓。上述具体证明除自反性相关部分外不使用 $\mathcal X$ 完备，题设 Banach 条件统一了本组习题的背景。
\end{remark}
% END SUPPLEMENT: id=v14-ch05-ex23-solution
% SOURCE p=175; printed=160; object=ch05-ex24-reflexive-dual; number=24; region=lower
\begin{exercise}\label{exercise:ch05-reflexive-dual}
设 $\mathcal X$ 为 Banach 空间。
\begin{enumerate}
\item\label{item:ch05-triple-dual-splitting}
\begin{definition*}[本题的自然像及其零化子]\phantomsection\label{definition:ch05-triple-dual-annihilator}
以 $\widehat{\mathcal X},\widehat{\mathcal X^*}$ 分别表示 $\mathcal X,\mathcal X^*$ 在 $\mathcal X^{**},\mathcal X^{***}$ 中的自然像，并令 $\widehat{\mathcal X}^{\,0}=\{F\in\mathcal X^{***}:F|_{\widehat{\mathcal X}}=0\}$。
\end{definition*}
则 $\widehat{\mathcal X^*}\cap\widehat{\mathcal X}^{\,0}=\{0\}$，且 $\widehat{\mathcal X^*}+\widehat{\mathcal X}^{\,0}=\mathcal X^{***}$。
\item\label{item:ch05-reflexive-iff-dual} $\mathcal X$ 自反当且仅当 $\mathcal X^*$ 自反。
\end{enumerate}
\end{exercise}
\Needspace{8\baselineskip}
% SUPPLEMENT: id=v14-ch05-ex24-solution; parent=ch05-ex24-reflexive-dual; kind=solution+motivation+hypotheses; origin=independent-derivation
\begin{solution}
\textbf{逐项补写（重构新增）。}
\begin{enumerate}
\item 对应\subcref{exercise:ch05-reflexive-dual}{item:ch05-triple-dual-splitting}：对任意 $F\in\mathcal X^{***}$ 定义 $f(x)=F(\widehat x)$。它线性且 $|f(x)|\le\|F\|\|x\|$，所以 $f\in\mathcal X^*$。其自然像 $\widehat f\in\mathcal X^{***}$ 满足 $\widehat f(G)=G(f)$，其中 $G\in\mathcal X^{**}$。因而
\begin{align*}
(F-\widehat f)(\widehat x)=F(\widehat x)-\widehat x(f)=f(x)-f(x)=0.
\end{align*}
故 $F=\widehat f+(F-\widehat f)$ 给所需和分解。若 $\widehat f$ 同时在 $\widehat{\mathcal X}^{\,0}$ 中，则每个 $x$ 都满足 $f(x)=\widehat f(\widehat x)=0$，故 $f=0$、$\widehat f=0$，交只有零元。
\item 对应\subcref{exercise:ch05-reflexive-dual}{item:ch05-reflexive-iff-dual}：若 $\mathcal X$ 自反，则 $\widehat{\mathcal X}=\mathcal X^{**}$，其零化子为零。第一问的和分解即给 $\mathcal X^{***}=\widehat{\mathcal X^*}$，所以 $\mathcal X^*$ 自反。

反过来，若 $\mathcal X^*$ 自反，则 $\mathcal X^{***}=\widehat{\mathcal X^*}$。第一问的交为零从而给 $\widehat{\mathcal X}^{\,0}=\{0\}$。由于 $\mathcal X$ 完备，自然像 $\widehat{\mathcal X}$ 在 $\mathcal X^{**}$ 中闭。若它是真子空间，Hahn--Banach 分离就给一个非零 $F\in\mathcal X^{***}$ 消去它，与零化子为零矛盾。因此自然像为整个双对偶，$\mathcal X$ 自反。
\end{enumerate}
\end{solution}
\begin{remark}
\textbf{动机与条件（重构新增）。}三重对偶中任意泛函都有一部分由它在原空间自然像上的限制确定，另一部分完全消去该像。自反性正好排除第二部分。第一问不使用完备性；第二问逆向用 Banach 条件使自然像闭，才能把“没有非零零化子”加强为“自然像等于全空间”，否则该论证只能给稠密。
\end{remark}
% END SUPPLEMENT: id=v14-ch05-ex24-solution
% SOURCE p=175; printed=160; object=ch05-ex25-separable-dual; number=25; region=lower
\begin{exercise}\label{exercise:ch05-separable-dual}
若 $\mathcal X$ 为 Banach 空间且 $\mathcal X^*$ 可分，则 $\mathcal X$ 可分。（令 $\{f_n\}_{n=1}^\infty$ 为 $\mathcal X^*$ 的可数稠密子集；对每个 $n$ 选 $\|x_n\|=1$ 且 $|f_n(x_n)|\ge\frac12\|f_n\|$ 的 $x_n$。则 $\{x_n\}_{n=1}^\infty$ 的线性组合在 $\mathcal X$ 中稠密。）注意：$\mathcal X$ 可分并不蕴含 $\mathcal X^*$ 可分。
\end{exercise}
% SUPPLEMENT: id=v14-ch05-ex25-solution; parent=ch05-ex25-separable-dual; kind=solution+counterexample+motivation+hypotheses; origin=independent-derivation
\begin{solution}
\textbf{主断言与逆向反例补写（重构新增）。}若 $\mathcal X=\{0\}$，可分性显然。否则，当 $f_n\ne0$ 时按范数上确界选题设的单位向量；$f_n=0$ 时任取单位向量。令 $\mathcal M=\overline{\operatorname{span}\{x_n:n\ge1\}}$。若 $\mathcal M\ne\mathcal X$，分离定理给一个非零 $f\in\mathcal X^*$ 消去 $\mathcal M$。选 $n$ 使 $\|f_n-f\|<\|f\|/4$，则
\begin{align*}
\tfrac12\|f_n\|\le|f_n(x_n)|=|(f_n-f)(x_n)|<\tfrac14\|f\|,
\end{align*}
但 $\|f_n\|\ge\|f\|-\|f_n-f\|>3\|f\|/4$，所以左边大于 $3\|f\|/8$，矛盾。因此 $\mathcal M=\mathcal X$。

取这些 $x_n$ 的有限线性组合，系数限于 $\mathbb Q$（实域）或 $\mathbb Q+i\mathbb Q$（复域），得到可数集 $D$。任意有限标量组合可逐系数作有理逼近；因 $\|\sum_j(a_j-b_j)x_j\|\le\sum_j|a_j-b_j|$，可使范数误差任意小。故 $D$ 在整个 $\mathcal X$ 中稠密。

逆向以 $\mathcal X=\ell^1$ 为例。它是自然数集上计数测度的 $L^1$，因而为 Banach 空间；有理实部、虚部的有限支集序列构成可数稠密集，因为先截去绝对可和尾部，再逼近有限个坐标。

对每个有界标量序列 $a=(a_n)$，定义 $f_a(x)=\sum_{n=1}^\infty a_nx_n$。级数绝对收敛且 $|f_a(x)|\le\|a\|_\infty\|x\|_1$，所以 $f_a\in(\ell^1)^*$。测试单位坐标向量得 $\|f_a\|\ge|a_n|$，再取上确界得 $\|f_a\|=\|a\|_\infty$。若 $a\ne b$ 是两个只取0、1的序列，就有 $\|f_a-f_b\|=1$。此类序列不可数（若可数列举，逐坐标取反的对角序列不在列表中）。

可分度量空间不能含不可数个彼此距离至少1的点：可数稠密集的半径 $1/3$ 开球覆盖全空间，每个球至多含一个这样的点。因此 $(\ell^1)^*$ 不可分。这已充分证明题末的否定断言，无需预先使用整个对偶与 $\ell^\infty$ 的表示定理。
\end{solution}
\begin{remark}
\textbf{动机与条件（重构新增）。}给每个稠密测量挑一个近乎最大读数的单位向量，若这些向量没有稠密生成全空间，就能找到漏掉的非零测量，与测量族稠密矛盾。这里是对偶的\emph{范数}可分性；较弱拓扑的可分性不足以给上述统一误差界。证明主断言实际不使用 $\mathcal X$ 完备。零空间需单列，因为题目提示中的单位向量在其中不存在。逆向反例说明向量可数逼近不意味着所有有界线性测量也可数逼近。
\end{remark}
% END SUPPLEMENT: id=v14-ch05-ex25-solution
\Needspace{12\baselineskip}
% SOURCE p=175; printed=160; object=ch05-ex26-krein-extension; number=26; region=bottom
\begin{exercise}\label{exercise:ch05-krein-extension}
设 $\mathcal X$ 为实向量空间，$P\subseteq\mathcal X$ 满足：(i) $x,y\in P$ 蕴含 $x+y\in P$；(ii) $x\in P$ 且 $\lambda\ge0$ 蕴含 $\lambda x\in P$；(iii) $x\in P$ 且 $-x\in P$ 蕴含 $x=0$。
% SOURCE p=176; printed=161; object=ch05-ex26-krein-extension-continuation; region=top
（例如，若 $\mathcal X$ 是实值函数空间，可取 $P$ 为其中的非负函数集合。）
\begin{enumerate}
\item\label{item:ch05-cone-partial-order} 由 $x\le y$ 当且仅当 $y-x\in P$ 定义的关系是 $\mathcal X$ 上的偏序。
\item\label{item:ch05-krein-positive-extension} \textbf{Krein 延拓定理。}设 $\mathcal M$ 为 $\mathcal X$ 的子空间，并且对每个 $x\in\mathcal X$ 都存在 $y\in\mathcal M$ 使 $x\le y$。若 $f$ 为 $\mathcal M$ 上的线性泛函，且 $x\in\mathcal M\cap P$ 时 $f(x)\ge0$，则存在 $\mathcal X$ 上的线性泛函 $F$，使 $x\in P$ 时 $F(x)\ge0$，并且 $F|_{\mathcal M}=f$。（考虑 $p(x)=\inf\{f(y):y\in\mathcal M\text{ 且 }x\le y\}$。）
\end{enumerate}
\end{exercise}
% SUPPLEMENT: id=v14-ch05-ex26-solution; parent=ch05-ex26-krein-extension; kind=solution+source-correction+motivation+hypotheses; origin=independent-derivation
\begin{remark}
\textbf{空集边界的修正（重构新增）。}原题前三个条件都被 $P=\varnothing$ 真空满足，但此时没有 $x\le x$，第一问不是偏序。第一问需补 $P\ne\varnothing$，等价于在所列条件下补 $0\in P$。第二问的支配条件已经排除 $P=\varnothing$，无需另添假设。下文先在修正后的非空情形证明第一问，再按第二问原有假设证明延拓。
\end{remark}
\begin{solution}
\textbf{逐项完整补写（重构新增）。}
\begin{enumerate}
\item 对应\subcref{exercise:ch05-krein-extension}{item:ch05-cone-partial-order}：取 $z\in P$，由(ii)取 $\lambda=0$ 得 $0\in P$，于是 $x-x\in P$，关系自反。若 $x\le y$ 且 $y\le z$，则 $(z-y)+(y-x)=z-x\in P$，所以传递。若 $x\le y$ 且 $y\le x$，则 $y-x$ 及其负向量均在 $P$，由(iii)得 $y-x=0$，所以反对称。三条给偏序。
\item 对应\subcref{exercise:ch05-krein-extension}{item:ch05-krein-positive-extension}：先验证提示的 $p$ 为有限实值。对每个 $x$，支配条件使集合 $A_x=\{y\in\mathcal M:x\le y\}$ 非空。对 $-x$ 再使用条件，取 $z\in\mathcal M$ 使 $-x\le z$，于是 $-z\le x$。对任意 $y\in A_x$，有 $y+z\in\mathcal M\cap P$，所以 $f(y)+f(z)\ge0$，即 $f(y)\ge-f(z)$。因此 $p(x)>-\infty$；取一个 $y\in A_x$ 又给有限上界，故 $p(x)\in\mathbb R$。

若 $x\in\mathcal M$，则 $x\in A_x$，所以 $p(x)\le f(x)$；对任意 $y\in A_x$，$y-x\in\mathcal M\cap P$，故 $f(y)\ge f(x)$。因此 $p(x)=f(x)$，特别 $p(0)=0$。

对 $\lambda>0$，$y\mapsto\lambda y$ 为 $A_x$ 到 $A_{\lambda x}$ 的双射，因为正数倍乘及其逆都保持 $P$。故 $p(\lambda x)=\lambda p(x)$；$\lambda=0$ 由 $p(0)=0$ 成立。任给 $\varepsilon>0$，选 $y\in A_x$、$z\in A_w$，使 $f(y)<p(x)+\varepsilon$、$f(z)<p(w)+\varepsilon$。则 $y+z\in A_{x+w}$，因而
\begin{align*}
p(x+w)\le f(y+z)=f(y)+f(z)<p(x)+p(w)+2\varepsilon.
\end{align*}
让 $\varepsilon\downarrow0$ 得次可加性。所以 $p$ 为次线性泛函，且 $f=p$ 于 $\mathcal M$。

由实 Hahn--Banach 定理，存在实线性延拓 $F$，使 $F\le p$。若 $x\in P$，则 $-x\le0$，所以 $0\in A_{-x}$，进而 $p(-x)\le f(0)=0$。因此 $-F(x)=F(-x)\le p(-x)\le0$，即 $F(x)\ge0$。这证明所需正性及延拓结论。
\end{enumerate}
\end{solution}
\begin{remark}
\textbf{动机与逐条件说明（重构新增）。}加法封闭性产生传递性，非负缩放封闭性使正负方向的大小可以比较，尖性(iii)用于反对称；非空性补足自反性。子空间 $\mathcal M$ 的支配条件既给每个 $A_x$ 一个元素，也通过对 $-x$ 的支配给其泛函值下界，保证下确界是实数而非无穷。$f$ 的正性使这些上、下界相容；线性性使 $p$ 在 $\mathcal M$ 上回到 $f$。延拓正性最后通过控制 $-x$ 得到，不能把 $F(x)\le p(x)$ 直接误读成正性。第二问的代数证明实际上无需尖性，尖性用于将背景关系称为偏序。整个论证不涉及范数、拓扑或完备性。
\end{remark}
% END SUPPLEMENT: id=v14-ch05-ex26-solution
\end{document}
